% ===================================================================
%  TABULKA č. 15 — Trh. — Potraviny.
%  Traduction 2026 d'apres texto/io, controlee sur texto/fr, texto/ru
%  et texto/uk. Consigne : voir texto/cs/10-tabulka-01.tex.
%
%  LE TABLEAU DU MARCHE EST CELUI OU LA REGLE DU LIEU D'APPRENTISSAGE
%  SE VERIFIE LE MIEUX DE TOUTE LA COLONNE, comme elle s'etait
%  verifiee au meme tableau en irlandais et en galicien : la
%  nourriture s'apprend a table. Cent objets, et l'emprunt s'arrete
%  aux deux frontieres de l'assiette — le the, le chocolat, le cafe,
%  les bananes, l'ananas d'un cote ; tout le reste est tcheque.
%
%  ET LE TCHEQUE AJOUTE ICI CE QUE LE GALICIEN AVAIT MONTRE SUR
%  L'ETAL DU POISSONNIER, mais a l'autre bout de la halle. L'ido y
%  distingue CINQ produits de charcuterie — « shinko », « sango-
%  sociso », « sociseto », « trip-sociseto », « lardo » — et le
%  tcheque a pour chacun un mot herite et distinct : « šunka »,
%  « jelito », « klobása », « jitrnice », « slanina ». Aucune des
%  douze colonnes precedentes n'y arrive sans emprunt ni
%  approximation ; « jitrnice », le boudin de foie et de gruau, n'a
%  d'equivalent simple nulle part ailleurs dans le volume.
%
%  LA REGLE EST CELLE DE LA VENDANGE AU TABLEAU 8 : la ou l'on vit de
%  la chose, on distingue. La Boheme vit du cochon, non de la mer.
% ===================================================================

%%K t15-apar-1 apar
\VUsaut{4.0mm}
\VUtitre{15.0pt}{60}{TABULKA č. 15}
\VUsaut{3.0mm}

%%K t15-tit-1 sub
\VUpk{11.4pt}{40}{Trh. --- Potraviny.}
\VUsaut{3.0mm}

%%K t15-01-1 p
1. --- Většina obchodníků, kteří nám dodávají zboží potřebné k naší výživě, se shromažďuje pod
\VUgras{halou} \textsuperscript{(1)} \VUgras{krytého trhu} nebo v sousedních ulicích.

%%K t15-02-1 p
2. --- \VUgras{Uzenář} \textsuperscript{(2)}, který prodává \VUgras{šunku}
\textsuperscript{(3)}, \VUgras{jelito} \textsuperscript{(4)}, \VUgras{klobásu}
\textsuperscript{(5)}, \VUgras{jitrnici} \textsuperscript{(6)} a \VUgras{slaninu}
\textsuperscript{(7)}, stojí blízko \VUgras{prodavačky másla a sýra} \textsuperscript{(8)},
jejíž pult je zásoben \VUgras{holandskými sýry} \textsuperscript{(9)} s červenou kůrkou,
\VUgras{camembertskými sýry} \textsuperscript{(10)}, \VUgras{bochníky ementálu}
\textsuperscript{(11)} a čerstvými
\VUgras{hroudami másla} \textsuperscript{(12)}.

%%K t15-03-1 p
3. --- Před ní nabízí \VUgras{prodavačka zeleniny} \textsuperscript{(13)} svým zákaznicím
krásná \VUgras{rajčata} \textsuperscript{(14)}, \VUgras{květáky} \textsuperscript{(15)},
\VUgras{salát} \textsuperscript{(16)}, \VUgras{hlávkové zelí} \textsuperscript{(17)},
\VUgras{tykvičky} \textsuperscript{(19)}, \VUgras{melouny} \textsuperscript{(18)} a dokonce
\VUgras{houby} \textsuperscript{(20)}. Ale
\VUgras{pivovarský podomek}, který zastavil svůj \VUgras{rozvážkový vůz}
\textsuperscript{(21)} naložený \VUgras{sudy piva} \textsuperscript{(22)} právě před jejím
pultem, brání jejím zákazníkům přiblížit se. Zahradnice jí přináší velký koš
\VUgras{artyčoků} \textsuperscript{(23)}.

%%K t15-tit-2 sub
\VUsaut{4.0mm}
\VUpk{10.2pt}{40}{Prodej a nákup.}
\VUsaut{3.0mm}

%%K t15-04-1 p
4. --- Na trhu je velký ruch, protože přišla hodina \VUgras{dražby} \textsuperscript{(24)}.
\VUgras{Hospodyně} \textsuperscript{(26)}, které sem přicházejí nakupovat, se zastavují cestou
před krámem \VUgras{prodavačky drštěk} \textsuperscript{(25)}, aby koupily \VUgras{dršťky}
nebo \VUgras{telecí hlavu.}

%%K t15-05-1 p
5. --- Tlustý \VUgras{kuchař} \textsuperscript{(27)} smlouvá o velmi krásného
\VUgras{tuňáka} s \VUgras{prodavačkou ryb} \textsuperscript{(28)}, která podle své
libovůle zvyšuje své ceny. Dostala velké množství \VUgras{úhořů, mořských jazyků, pstruhů} a
\VUgras{kaprů.} Její \VUgras{treska} a její \VUgras{slávky} jsou velmi čerstvé.

%%K t15-tit-3 sub
\VUsaut{4.0mm}
\VUpk{10.2pt}{40}{Laciné a drahé zboží.}
\VUsaut{3.0mm}

%%K t15-06-1 p
6. --- \VUgras{Prodavačka drůbeže} \textsuperscript{(29)}, usazená blízko ní, je velmi
svědomitá. Když prodává \VUgras{krocana, husu, kachnu} nebo prosté \VUgras{kuře,} žádá jen
spravedlivou cenu.

%%K t15-07-1 p
7. --- Na rohu náměstí, v prvním patře, se před nedávnem usadil \VUgras{prodavač plátna}
\textsuperscript{(30)}, u něhož jsou kupující ženy vždy jisty, že najdou velmi krásný výběr
\VUgras{ubrusů,} \VUgras{ubrousků, zástěr} a \VUgras{utěrek.} V témž patře si \VUgras{nožíř}
\textsuperscript{(31)} zřídil svou dílnu. Brousí \VUgras{nože} prodavaček v hale a vyrábí pro
zahradníky \VUgras{zahradnické nože} z nejtvrdší
\VUgras{oceli.}

%%K t15-tit-4 sub
\VUsaut{4.0mm}
\VUpk{10.2pt}{40}{Koloniál.}
\VUsaut{3.0mm}

%%K t15-08-1 p
8. --- Pod jeho dílnou byl před nedávnem otevřen velký obchod s potravinami. Už to není prostý
a bezvýznamný \VUgras{koloniál} \textsuperscript{(32)} jako kdysi, který prodával jen běžné
zboží, \VUgras{čaj} \textsuperscript{(33)}, \VUgras{čokoládu} \textsuperscript{(34)},
\VUgras{homolový cukr} \textsuperscript{(37)} a \VUgras{mýdlo} \textsuperscript{(38)}. Tam se
prodává jakékoli zboží, \VUgras{křupavé sušenky},
\VUgras{konzervy} \textsuperscript{(35)}, \VUgras{likéry} \textsuperscript{(36)},
\VUgras{rum,} \VUgras{koňak, kirš, anýzovka} a \VUgras{curaçao.} Tam se najde zvěřina:
\VUgras{zajíc} \textsuperscript{(39)}, \VUgras{drozdi} \textsuperscript{(40)},
\VUgras{koroptve} \textsuperscript{(41)}, \VUgras{křepelky} \textsuperscript{(42)},
\VUgras{divoká kachna} \textsuperscript{(43)} nebo \VUgras{bažant}
\textsuperscript{(44)}, právě tak snadno jako exotické ovoce jako \VUgras{banány}
\textsuperscript{(46)} a \VUgras{ananasy.}

%%K t15-09-1 p
9. --- Velmi ochotný \VUgras{příručí} koloniálu \textsuperscript{(45)} je připraven dát, co si
kdo přeje: \VUgras{zavařeninu} \textsuperscript{(47)} rybízovou nebo kdoulovou, \VUgras{sádlo}
\textsuperscript{(48)}, \VUgras{okurky} \textsuperscript{(49)} nebo solené \VUgras{sardinky}
\textsuperscript{(50)}.

%%K t15-09-2 p
To je velmi příhodné pro \VUgras{zákaznice} \textsuperscript{(51)}, které vedle nejběžnějšího
zboží najdou nejvzácnější rané plody a nejchutnější ovoce: šťavnaté
\VUgras{hrušky} \textsuperscript{(52)}, \VUgras{švestky} \textsuperscript{(53)},
\VUgras{hrozny} \textsuperscript{(54)}, \VUgras{broskve} \textsuperscript{(55)},
\VUgras{mandle} \textsuperscript{(56)} a \VUgras{ořechy} \textsuperscript{(57)}.

%%K t15-tit-5 sub
\VUsaut{4.0mm}
\VUpk{10.2pt}{40}{Zákazníci.}
\VUsaut{3.0mm}

%%K t15-10-1 p
10. --- \VUgras{Rýže} \textsuperscript{(58)} a \VUgras{čočka} \textsuperscript{(59)} jsou
vedle bedny uzených \VUgras{sleďů} \textsuperscript{(60)}; \VUgras{brambor}
\textsuperscript{(61)} a \VUgras{fazole} \textsuperscript{(63)} sousedí s
\VUgras{jablky} \textsuperscript{(62)}, \VUgras{fíky} \textsuperscript{(64)},
\VUgras{sušenými švestkami} \textsuperscript{(65)} a \VUgras{lískovými oříšky}
\textsuperscript{(66)}. V tom obchodě se najdou čerstvá \VUgras{vejce} \textsuperscript{(67)}
a \VUgras{olivy} \textsuperscript{(68)}; proto se \VUgras{koše} \textsuperscript{(70)} a
\VUgras{nákupní síťky} \textsuperscript{(73)} naplňují velmi rychle.

%%K t15-11-1 p
11. --- \VUgras{Káva} \textsuperscript{(72)} té výborné firmy získala mezi
\VUgras{služkami} \textsuperscript{(69)} a \VUgras{kuchařkami} zaslouženou pověst,
protože starý \VUgras{koloniálník} \textsuperscript{(71)} ji sám co nejpečlivěji praží.

%%K t15-tit-6 sub
\VUsaut{4.0mm}
\VUpk{10.2pt}{40}{Podívané.}
\VUsaut{3.0mm}

%%K t15-12-1 p
12. --- Ne všichni jdou na trh, aby se zásobili. V malebném ruchu uličky, která vede k hale,
je vidět nejrozmanitější osoby. Vedle vlídného \VUgras{řeznického pomocníka}
\textsuperscript{(75)}, který před sebou tlačí \VUgras{kolečko} \textsuperscript{(74)}, jdou
dva vojáci na dovolené, velmi hrdí na to, že mohou ukázat svou skvělou uniformu:
\VUgras{husar} \textsuperscript{(76)} s modrým
\VUgras{dolmanem} a \VUgras{čákou s chocholem,} a \VUgras{zuávský desátník,} s
nabíranými rukávy a s hlavou pokrytou \VUgras{fezem.}

%%K t15-13-1 p
13. --- \VUgras{Pekař} \textsuperscript{(80)}, který stojí na prahu \VUgras{pekárny}
\textsuperscript{(78)}, právě uhnětl těsto svého \VUgras{chleba} \textsuperscript{(79)} a
vyňal z pece \VUgras{rohlíky} \textsuperscript{(81)}, \VUgras{housky} \textsuperscript{(82)} a
\VUgras{kulaté bochníky} \textsuperscript{(83)}, které jsou ve výkladu.

%%K t15-14-1 p
14. --- Nad pekárnou bydlí obratná \VUgras{švadlena prádla} \textsuperscript{(84)}, která
zaměstnává několik \VUgras{vyšívaček} \textsuperscript{(85)}. Vedle ní bydlí
\VUgras{pradlena a žehlířka} \textsuperscript{(86)}, která škrobí \VUgras{prádlo}
\textsuperscript{(88)} po vyprání a usušení a která je hladí \VUgras{žehličkou}
\textsuperscript{(87)}.

%%K t15-tit-7 sub
\VUsaut{4.0mm}
\VUpk{10.2pt}{40}{Maso, ryby a zelenina.}
\VUsaut{3.0mm}

%%K t15-15-1 p
15. --- V \VUgras{řeznictví} \textsuperscript{(89)}, umístěném v přízemí, se prodává všeliké
\VUgras{maso, skopové} \textsuperscript{(90)}, \VUgras{hovězí} \textsuperscript{(94)} nebo
\VUgras{telecí} \textsuperscript{(92)} v podobě
\VUgras{skopových kýt, bifteků, pečení} a \VUgras{kotlet.} \VUgras{Řeznický pomocník}
\textsuperscript{(93)} rozřezává všechna ta masa a klade stranou \VUgras{morkové kosti}
\textsuperscript{(96)}. U dveří řeznictví je \VUgras{prodavačka ústřic}
\textsuperscript{(97)}, která prodává \VUgras{ústřice} \textsuperscript{(98)}. Otvírá je,
přeje-li si to kdo.

%%K t15-16-1 p
16. --- Všichni ti obchodníci platí koncesi nebo poplatek za místo; proto
\VUgras{strážník} \textsuperscript{(100)} nemilosrdně stíhá ubohé \VUgras{pouliční
prodavačky zeleniny} \textsuperscript{(99)}, které jim konkurují rozvážejíce na svých
vozíčcích \VUgras{mrkev, ředkvičky, tuřín, pórek} nebo
\VUgras{svazky chřestu, čerstvý hrášek, špenát, šťovík, řeřichu} a \VUgras{petržel.}

%%K t15-tit-8 sub
\VUsaut{4.0mm}
\VUpk{10.2pt}{40}{Pozor na strážníka!}
\VUsaut{3.0mm}

%%K t15-17-1 p
17. --- Chlapec, který prodává \VUgras{česnek} \textsuperscript{(101)} a
\VUgras{cibuli,} velmi dobře ví, že ani on nesmí prodávat své zboží blízko trhu. Utíká
a při běhu riskuje, že převrhne koš \VUgras{krabů,} \VUgras{raků} a \VUgras{krevet}
\VUgras{prodavače} \textsuperscript{(102)} \VUgras{langust} a \VUgras{humrů.}

%%K t15-18-1 p
18. --- Všichni se hemží a hlučí; ale to nebrání zřetelně slyšet hlas potulného
\VUgras{skláře} \textsuperscript{(103)} ani volání \VUgras{uhlíře}
\textsuperscript{(104)}, který právě na několik okamžiků opustil svůj vozík, aby dodal
\VUgras{pytel uhlí} \textsuperscript{(105)}.
